\section{Layers and positive divisors}
\label{sec:layers}

We develop the layer calculus needed to compare representatives of a
parabolic coset.  The underlying linear heart and twist calculations
are related to \cite[Sections~5 and~6.1]{HengLicata2024}.
We prove the statements used here, with particular attention to the degree conventions
and to the arbitrary inputs allowed in the cohomological divisor test.
None of the results in this section requires the initial projective to
be attached to every operated color.

\subsection{The linear heart}

Work in the homotopy category of bounded complexes of finite graded left
projective $Z$-modules.  A differential has cochain degree $1$ and
internal degree $0$.  The shift $\langle u\rangle$ raises internal
degrees by $u$, and $[1]$ lowers cochain placements by one.  A generator
of $P_\alpha\langle u\rangle$ in cochain placement $c$ has
\emph{level}
\begin{equation}
 d=c+u.\label{lay:level}
\end{equation}
Throughout, generator dimensions mean multiplicities of projective
summands, not dimensions of those projectives as vector spaces.

A complex is \emph{minimal} if its differential entries belong to the
graded radical.  Gaussian elimination removes an invertible entry
between shifted indecomposable projectives by splitting off a
contractible two-term complex.  Since the number of summands is finite,
this produces a minimal model.  A homotopy equivalence between minimal
models is an isomorphism modulo the radical in every cochain degree:
the homotopy terms in its composites with a homotopy inverse contain a
radical differential.  Consequently its term maps are invertible.
Thus a minimal model is unique up to chain isomorphism, in particular
with uniquely determined generator multiplicities and placements.

In a minimal model the differential has only path lengths $1$ and $2$.
An entry of length $a$ changes generator placements by $(1,-a)$, hence
changes the level by $1-a$.  The linear differential preserves level,
and the token differential lowers it by one.  Define $Y^d(E)$ to be the
generators of level $d$ with their linear differential.  The path-length
$2$ part of the equation $\partial^2=0$ says that this is a complex.

For clarity, its coefficient description has spaces
\[
 V^d_{\alpha,c}(E),
 \qquad
 P_\alpha\langle d-c\rangle\otimes V^d_{\alpha,c}(E)
 \text{ in cochain placement }c.
\]
Each arrow entry gives a coefficient map from
$V^d_{\alpha,c}(E)$ to $V^d_{\beta,c+1}(E)$.  At each vertex these
maps satisfy the backtrack relation obtained from the opposite-arrow
pairing in $Z^2$; all other products of two arrows in $Z$ vanish.
We call such a finite graded array a \emph{coefficient representation}.
Its arrow maps raise cochain placement and are therefore jointly
nilpotent.  Its top and socle at $\alpha$ are, respectively,
\[
 \Top_\alpha M=
 \frac{\bigoplus_c V_{\alpha,c}(M)}
      {\text{images of incoming coefficient maps}},
 \qquad
 \soc_\alpha M=
 \bigcap\ker(\text{outgoing coefficient maps}),
\]
with the inherited placements.  Write $\dim M$ for the vector of
generator multiplicities summed over placements.

\begin{lemma}[Layer structure]\label{lay:structure}
The full subcategories defined by minimal generator levels at most $0$
and at least $0$ form a bounded $t$-structure.  The truncation
$\tau_{\le k}E$ is the subcomplex of levels at most $k$, and the
complementary quotient is $\tau_{\ge k+1}E$.

The heart is the abelian category of finite graded coefficient
representations.  In particular, subobjects of a pure layer are exactly
its graded subrepresentations.  Its unshifted cohomology layer at $d$
is $Y^d(E)$.
\end{lemma}

\begin{proof}
Differentials do not raise level, so the indicated low terms form a
subcomplex.  A degree-zero map from a generator of level $d$ to one of
level $e$ requires a path of length $d-e$.  There are no such maps from
levels at most $0$ to levels at least $1$.  This proves the required
orthogonality.  The cochain shift $[1]$ lowers levels by one, giving the
shift-closure axioms, and the short exact sequence of low subcomplex,
full complex and high quotient gives the truncation triangle.  Finitely
many occupied levels give boundedness.

A degree-zero map between pure complexes at the same level has only
length-zero entries.  It is precisely a map of graded coefficient
representations.  A homotopy between two such maps would require path
length $-1$, so there are no homotopies to quotient out.  Conversely a
finite coefficient representation gives a pure complex, since its
backtrack relations say exactly that the differential squares to zero.
Kernels and cokernels, formed in each coefficient space, retain the
relations.  This identifies the heart, its subobjects and its
cohomology layers.
\end{proof}

We use the long exact cohomology sequence associated with a distinguished
triangle in a $t$-structure \cite[Section~1.3, Theorem~1.3.6]{BBD1982}.
Denote its heart-valued cohomology by
$H^d_{\mathrm{lay}}$; with our conventions,
\[
 H^d_{\mathrm{lay}}(E)=Y^d(E)[d],
 \qquad
 H^d_{\mathrm{lay}}(E[1])=H^{d+1}_{\mathrm{lay}}(E).
\]
An exact functor has \emph{layer amplitude} $[a,b]$ if it sends a pure
layer at $d$ into levels $[d+a,d+b]$.  Bounded truncation then gives the
corresponding bounds on any interval of input levels.

\begin{lemma}[Adjacent layers]\label{lay:adjacent}
For an exact functor $F$ of layer amplitude $[0,1]$ there is a short
exact sequence of pure level-$d$ objects
\begin{equation}
\begin{split}
0\longrightarrow Y^d\bigl(F(Y^{d-1}(E))\bigr)
\longrightarrow Y^d(FE)
\longrightarrow Y^d\bigl(F(Y^d(E))\bigr)
\longrightarrow0.
\end{split}\label{lay:adjacent-positive}
\end{equation}
For amplitude $[-1,0]$ the corresponding sequence is
\begin{equation}
\begin{split}
0\longrightarrow Y^d\bigl(F(Y^d(E))\bigr)
\longrightarrow Y^d(FE)
\longrightarrow Y^d\bigl(F(Y^{d+1}(E))\bigr)
\longrightarrow0.
\end{split}\label{lay:adjacent-negative}
\end{equation}
All terms are shifted by $[d]$ to regard these as sequences in the
heart.  In particular, generator dimensions are additive.
\end{lemma}

\begin{proof}
In the positive case apply $F$ to the triangle
\[
 \tau_{\le d-1}E\longrightarrow\tau_{\le d}E
 \longrightarrow Y^d(E).
\]
The first term after applying $F$ has levels at most $d$; the last has
levels $d,d+1$.  Hence the flanking groups
$H^{d-1}_{\mathrm{lay}}F(Y^d(E))$ and
$H^{d+1}_{\mathrm{lay}}F(\tau_{\le d-1}E)$ vanish, giving a short
exact sequence at $d$.  Input levels at most $d-2$ contribute nothing
there, identifying the left term in
\eqref{lay:adjacent-positive}; input levels at least $d+1$ contribute
nothing there, identifying its middle term.  For the negative case
apply the same argument to
$\tau_{\le d}E\to\tau_{\le d+1}E\to Y^{d+1}(E)$.
After $F$, its first term has levels at most $d$ and its last has levels
$d,d+1$.  Discarding input levels below $d$ and above $d+1$ gives
\eqref{lay:adjacent-negative}.
\end{proof}

\begin{lemma}[Composition at the bottom]\label{lay:bottom}
Let $F_1,\ldots,F_r$ be exact functors with no negative change of
level.  On input pure at $d$, the level-$d$ part of
$F_r\cdots F_1$ is obtained by retaining only level $d$ after every
operation.  Relations between these composite functors induce the same
relations between their bottom-layer operations.
\end{lemma}

\begin{proof}
After the first operation the quotient by the level-$d$ truncation
has levels at least $d+1$.  All subsequent operations preserve that
bound.  Its long exact layer sequence shows that discarding this
quotient does not change the final level $d$.  Repeat after each
operation.
\end{proof}

\subsection{Extremal tests and single twists}

For a bigraded vector complex $V$, write $H^t(V)$ for the direct sum
of its cohomology spaces with cochain degree plus internal degree
equal to $t$.  This \emph{total test degree} is distinct from the
cochain degree alone.  Differentials raise total degree by one; only
cochain degrees determine their signs.

\begin{lemma}[Extremal tests]\label{lay:extremal}
If $M$ is pure at level $d$, then
\begin{equation}
 H^d(e_\alpha M)\cong\soc_\alpha M,
 \qquad
 H^{d+2}(e_\alpha M)\cong(\Top_\alpha M)\langle2\rangle.
 \label{lay:extremal-formula}
\end{equation}
The second shift records the token's internal placement.  More
generally, if a minimal complex $E$ has no generators above $d$ at
$\alpha$ or its neighbors, its degree-$(d+2)$ test is the top of
$Y^d(E)$ at $\alpha$, shifted internally by $2$, and all higher test
cohomology there vanishes.  If there are no generators below $d$ at
those vertices, its degree-$d$ test is the socle of $Y^d(E)$ there.

A nonzero pure layer supported on a vertex set has a nonzero top and
a nonzero socle somewhere in that set.  The same assertions hold for
right projective complexes, tested by $Me_\alpha$.
\end{lemma}

\begin{proof}
The test of a pure layer has path lengths $0,1,2$, in total degrees
$d,d+1,d+2$.  Its length-zero cycles are the common kernels of the
outgoing coefficient maps.  At the token end, the images of the
incoming differential are exactly the incoming coefficient images,
by the perfect opposite-arrow pairings.  This proves
\eqref{lay:extremal-formula}.

Every term of $e_\alpha E$ comes from a generator at $\alpha$ or a
neighbor, and its total degree is the generator level plus the path
length.  Under the upper bound, only level-$d$ tokens can occur at
$d+2$.  A token differential into them from length zero would require
a generator at level $d+1$, which is absent.  The remaining boundaries
are the same incoming linear images.  The lower statement follows
from the length-zero end in the same way.  Finally, coefficient arrows
strictly raise cochain placement in a finite array.  A nonzero array
therefore cannot have zero socle everywhere, or be generated by its
incoming images everywhere.  Reversing module sides proves the right
statements.
\end{proof}

Let $F_i=B_{\sigma_i}(-)$ be the color twist of
Proposition~\ref{cat:action}.  Denote by $r_i$ the product of the commuting
coordinate reflections in its unfolded fiber.  Thus, on folded
vectors,
\[
 r_i\phi(v)=\phi\bigl(v-(Cv)_i\varepsilon_i\bigr).
\]

\begin{lemma}[Positive layers]\label{lay:positive}
Each $F_i$ has layer amplitude $[0,1]$.  For a pure input $M$ at $d$,
its upper layer is supported on color $i$, has zero linear
differential, and has multiplicity $\dim\Top_\alpha M$ at each
$\alpha$ of that color.  Its bottom layer has unchanged multiplicities
at all other vertices, and
\begin{equation}
 \dim Y^d(F_iM)_\alpha
 =\sum_\beta(2I-\widetilde C)_{\alpha\beta}\dim M_\beta
       -\dim M_\alpha+\dim\Top_\alpha M
 \quad(\alpha\text{ of color }i).
 \label{lay:positive-count}
\end{equation}
The bottom layer has zero socle in color $i$.  If $i\in J$ and $M$
has zero socle on every color of $J$, its bottom layer after $F_i$ has
zero socle on all of $J$.

For any bounded complex $E$ of graded left modules, with no
projectivity requirement on its terms,
\begin{equation}
 e_\alpha F_iE\simeq(e_\alpha E)\langle2\rangle[1]
 \quad(\alpha\text{ of color }i).
 \label{lay:positive-shift}
\end{equation}
There is an identical homotopy equivalence for right actions and right
tests.  Thus an acted test shifts up one in total degree and down one
in cochain degree.
\end{lemma}

\begin{proof}
For pure input at $d$, the evaluation cone inserts generators in the
acted color at levels $d+a-1$, where $a=0,1,2$ is the path length.
The length-zero insertions at $d-1$ map by identities, modulo the
radical, to the old generators in that color.  Cancelling those pairs
removes every generator below $d$.  No generator was inserted above
$d+1$, and only the acted color can occur above $d$.

For one vertex $\alpha$ of color $i$, put $A=e_\alpha E$.
Distinct vertices of that color are orthogonal, so the acted test is
\begin{equation}
 \Cone\bigl(A\oplus A\langle2\rangle
          \xrightarrow{(\mathrm{id},\theta_\alpha)} A\bigr).
 \label{lay:token-cone}
\end{equation}
Projection to the token source $A\langle2\rangle[1]$ is a chain map
with contractible kernel $\Cone(\mathrm{id}_A)$, proving
\eqref{lay:positive-shift}.  This is a homotopy calculation for any
$A$, so it neither assumes projectivity nor merely asserts a
quasi-isomorphism.  The right calculation uses the same cone.

For pure input, the degree-$(d+3)$ output test equals the old
degree-$(d+2)$ test.  Lemma~\ref{lay:extremal} identifies the latter
with the old top.  Since the upper layer is supported on a color with
no internal arrows, the former counts its generators.  This proves
the upper-layer assertion.

For the multiplicity formula, count generators with sign $(-1)^d$.
The signed contributions of a cancelled scalar pair sum to zero,
because its levels are consecutive.  In the evaluation cone the
inserted signed contribution from a generator at $\beta$ to an acted
vertex $\alpha$ is minus the alternating path count
\[
 2\delta_{\alpha\beta}-(2I-\widetilde C)_{\alpha\beta}.
\]
Consequently the total signed vector transforms by $r_i$.  On pure
input the bottom vector minus the upper vector is therefore
$r_i\dim M$, proving \eqref{lay:positive-count}.  The same calculation
also shows, for general inputs, that
\begin{equation}
 \sum_d(-1)^d\dim Y^d(F_iE)
       =r_i\sum_d(-1)^d\dim Y^d(E).
 \label{lay:signed-reflection}
\end{equation}

By \eqref{lay:positive-shift}, the degree-$d$ test at an acted vertex
is zero.  Lemma~\ref{lay:extremal} identifies it with the socle of
the bottom layer.  At another vertex $\beta$, the first term in the
tested evaluation triangle is a sum of
$(e_\alpha M)\langle1\rangle$, over adjacent acted vertices.
It has no degree-$d$ cohomology, and its degree-$(d+1)$ cohomology
is the sum of the old socles at those vertices.  If these socles and
the old socle at $\beta$ vanish, the long exact sequence gives zero
degree-$d$ cohomology after the move.  This proves preservation of
socle vanishing on $J$.
\end{proof}

In particular, Lemma~\ref{lay:adjacent} says that the new layer at $d$
under $F_i$ is just its bottom operation on $Y^d(E)$ whenever the
old top at color $i$ in level $d-1$ vanishes.  By
Lemma~\ref{lay:bottom} and the braid relations, the bottom operations
of positive twists themselves satisfy the positive braid relations.

\begin{lemma}[Inverse amplitude and exterior invariance]
\label{lay:inverse}
An exact equivalence of layer amplitude $[0,1]$ has inverse amplitude
$[-1,0]$.  In particular $F_i^{-1}$ has that amplitude, and
\begin{equation}
 e_\alpha F_i^{-1}E\simeq(e_\alpha E)\langle-2\rangle[-1]
 \quad(\alpha\text{ of color }i).
 \label{lay:negative-shift}
\end{equation}
This test identity is valid for arbitrary bounded module complexes
and also on the right.  For either sign, acting in color $i$ leaves
the minimal generator spaces at other colors unchanged up to
isomorphism, with both placements retained.  Formula
\eqref{lay:signed-reflection} also holds for $F_i^{-1}$.
\end{lemma}

\begin{proof}
Let $F$ be the equivalence and let $N$ lie in the heart.  For any $A$
in levels at most $-2$, its image $FA$ has levels at most $-1$, so
\[
 \Hom(A,F^{-1}N)=\Hom(FA,N)=0.
\]
The orthogonal characterization of the $t$-structure implies that
$F^{-1}N$ has levels at least $-1$.  If $B$ has levels at least $1$,
then so does $FB$, and
\[
 \Hom(F^{-1}N,B)=\Hom(N,FB)=0.
\]
Hence $F^{-1}N$ has levels at most $0$.  Shifting proves the amplitude
bound on every pure input.  Applying \eqref{lay:positive-shift} to
$F_i^{-1}E$ and using the bimodule homotopy inverse proves
\eqref{lay:negative-shift}, including for nonprojective inputs and
for right complexes.  Applying \eqref{lay:signed-reflection} and
$r_i^2=1$ proves the inverse reflection formula.

For a vertex $\beta$ outside color $i$, the right simple satisfies
\[
 U_\beta B_{\sigma_i}^{\pm1}\simeq U_\beta.
\]
Indeed the inserted projective bimodule terms have left vertex in
color $i$ and are annihilated by $U_\beta$.  Tensoring a minimal left
projective complex with $U_\beta$ kills its radical differential and
records precisely the generators at $\beta$ with both placements.
The displayed equivalence proves the exterior assertion.
\end{proof}

All assertions about projective layers have their right-handed
versions, obtained by using the opposite algebra.  In particular the
inverse layer-amplitude statement does not assert a projective-layer
description for a nonprojective input; for such an input the test
identities \eqref{lay:positive-shift} and
\eqref{lay:negative-shift} still hold.  The normalized layers containing
a residual right simple will be treated separately in
Section~\ref{sec:caps}.

\begin{lemma}[Compatibility with labels]\label{lay:labels}
Layer truncations commute with the exact label functors $T_H$.
For the family $E_L(x)$ of Proposition~\ref{cat:family}, the weighted generator,
top and socle dimensions in each layer are proportional on every
color fiber.  In particular, a top or socle occurring in one member
at a vertex of color $i$ occurs at every vertex of color $i$ in some
member of the family.  The same proportionality holds for families
constructed from this one by twists and layer truncations.
\end{lemma}

\begin{proof}
Tensoring a projective label by $H$ preserves its cochain and internal
placements.  It preserves radical path lengths, hence commutes with
the low-level subcomplex and the high-level quotient.  On coefficient
representations it is exact, since the label category is semisimple.
At a fixed color and placements, the coefficient spaces, tests and
their cohomology are obtained by applying
$\Hom_{\cC}(K,-)$ to the corresponding label objects.  Exact label
tensoring also tensors these objects and their cohomology.  The
weighted multiplicity identity
\[
 \sum_L w_L[D\otimes L:K]=w_K\dim_{\cC}D
\]
therefore gives proportionality for generators and for the extremal
tests, hence for tops and socles by Lemma~\ref{lay:extremal}.
All weights are positive, proving the occurrence assertion.
Twists and truncations commute with label tensoring, so the argument
applies to the constructed families as well.
\end{proof}

\subsection{The highest-test divisor}

We next recover a positive divisor from a highest cohomological test.
This is the mechanism used by Brav and Thomas to detect finite-type
Garside factors \cite[Proposition~3.1 and Lemma~3.3]{BravThomas2011}.
Here we prove the version for arbitrary bounded test inputs; right simple
modules in Section~\ref{sec:caps} will require precisely that generality.

Fix a spherical color set $J$.  By Lemma~\ref{cat:unfolding}, its
unfolded graph is finite simply laced.  In the next two lemmas the
letters are individual unfolded vertices; $B_\alpha$ denotes their
single-vertex evaluation twist.  For an input $E$ define
\begin{equation}
 h_J(E)=\max\{t:H^t(e_\alpha E)\ne0
                      \text{ for some }\alpha\in\widetilde J\},
 \label{lay:test-height}
\end{equation}
with value $-\infty$ when all these tests vanish.  For right inputs
use $Ee_\alpha$.  The following arguments require only bounded test
cohomology and the indicated action, not projective terms.  They apply
in particular to bounded complexes containing right simples.

\begin{lemma}[Change of test height]\label{lay:height-change}
If $h_J(E)$ is finite, then
\[
 h_J(E)\le h_J(B_\alpha E)\le h_J(E)+1.
\]
Height increases exactly when $\alpha$ was a highest-test vertex;
in that case it is the only new highest-test vertex.  Vanishing of all
tests on $\widetilde J$ persists.  The same statements hold for right
actions.  Consequently height is nondecreasing under positive color
words in $J$.
\end{lemma}

\begin{proof}
The acted test shifts up by one in total degree, and nonneighboring
tests are unchanged.  At a neighboring vertex $\beta$ the evaluation
triangle is
\begin{equation}
 (e_\alpha E)\langle1\rangle\longrightarrow e_\beta E
       \longrightarrow e_\beta B_\alpha E.
 \label{lay:neighbor-triangle}
\end{equation}
If the old maximum is $t$, the last term has no test cohomology above
$t$.  Thus precisely the asserted condition produces height $t+1$,
and only at $\alpha$.  If the old test at $\alpha$ has maximum $t-1$,
its shifted test reaches $t$.  If it has maximum below $t-1$, an old
degree-$t$ test elsewhere survives: in
\eqref{lay:neighbor-triangle} the map into that degree has zero source.
This proves monotonicity.  The same triangles preserve total
vanishing.  Reversing sides proves the right case; a color twist is
a product of the commuting unfolded twists in its fiber.
\end{proof}

\begin{lemma}[Top-divisor test]\label{lay:top-divisor}
Let $b$ be a positive word in the unfolded finite-type generators.
If $h_J(B_bE)>h_J(E)$, every vertex attaining the highest test
cohomology of $B_bE$ is a left divisor of $b$ in the positive Artin
monoid.  For a right action, such a vertex is a right divisor of $b$.
These conclusions hold for arbitrary inputs with bounded test
cohomology, including right-simple complexes.
\end{lemma}

\begin{proof}
We prove the left statement by induction on word length.  Write
$b=\alpha b'$ and $N=B_{b'}E$.  If the final action raises height,
its unique new highest vertex is $\alpha$ by
Lemma~\ref{lay:height-change}, and there is nothing to prove.
Otherwise
\[
 h_J(N)=h_J(B_\alpha N)=t>h_J(E).
\]
The vertex $\alpha$, when it is a highest vertex, already divides
$b$.  If a highest vertex $\beta\ne\alpha$ is not adjacent to
$\alpha$, it was a highest vertex of $N$.  By induction it divides
$b'$, and it commutes past $\alpha$.

Suppose instead that $\beta$ is adjacent to $\alpha$.  Since height
did not increase, $H^t(e_\alpha N)=0$.  The long exact sequence of
\eqref{lay:neighbor-triangle} gives
\begin{equation}
 H^t(e_\beta B_\alpha N)
 \cong\operatorname{coker}\bigl(
 H^{t-1}(e_\alpha N)\longrightarrow H^t(e_\beta N)\bigr).
 \label{lay:divisor-cokernel}
\end{equation}
This cokernel is nonzero, so $\beta$ is a highest vertex of $N$.
Induction gives a positive factorization $b'=\beta v$.  Put
$N_0=B_vE$, so $N\simeq B_\beta N_0$.

We now identify the arrow in \eqref{lay:divisor-cokernel}; this is the
step that permits nonprojective inputs.  Set
$A_0=e_\beta N_0$ and $C_0=e_\alpha N_0$, and choose opposite arrows
\[
 r\in e_\alpha Ze_\beta,\qquad
 s\in e_\beta Ze_\alpha,\qquad
 sr=\kappa\theta_\beta,\quad\kappa\ne0.
\]
The evaluated cones are
\[
\begin{split}
 e_\beta N&=\Cone\bigl(A_0\oplus A_0\langle2\rangle
                   \xrightarrow{(\mathrm{id},\theta_\beta)}A_0\bigr),\\
 e_\alpha N&=\Cone\bigl(A_0\langle1\rangle
                   \xrightarrow{r}C_0\bigr).
\end{split}
\]
Let $\pi:e_\beta N\to A_0\langle2\rangle[1]$ be the token
projection in \eqref{lay:token-cone}.  Multiplication by $s$ gives
$(e_\alpha N)\langle1\rangle\to e_\beta N$.  On its cone source,
the product $sr$ is $\kappa$ times the token, whereas its cone target
lands in the target $A_0$, killed by $\pi$.  Therefore $\pi s$ is,
up to the cochain sign and the nonzero scalar $\kappa$, exactly the
connecting projection of the second cone.

Since $\pi$ is a homotopy equivalence, the image of the map in
\eqref{lay:divisor-cokernel}, identified in
$H^{t-1}(e_\beta N_0)$, is
\[
 \ker\bigl(H^{t-1}(e_\beta N_0)
                 \longrightarrow H^t(e_\alpha N_0)\bigr).
\]
Its cokernel is nonzero, so this last arrow has nonzero image.
In particular $H^t(e_\alpha N_0)\ne0$.  By height monotonicity,
$h_J(N_0)\le h_J(N)=t$, hence
$h_J(N_0)=t>h_J(E)$.  Applying induction to $v$ gives
$v=\alpha v'$ in the positive monoid.  The simply-laced braid relation
now extracts the required vertex:
\[
 b=\alpha\beta\alpha v'=\beta\alpha\beta v'.
\]
This completes the induction.  On the right, use the right evaluation
cones and opposite arrows, with the final action on the right.  The
same proof gives right divisibility.
\end{proof}

\begin{remark}
Divisibility here is on the \emph{last-action side}: the left side for
left complexes and the right side for right complexes.  Passing from
an unfolded vertex divisor to an original color divisor requires the
finite-type greedy argument of Lemma~\ref{two:greedy}; preservation of
descents for reduced words alone does not suffice for an arbitrary
positive word.
\end{remark}

\begin{corollary}[Reduced-word amplitude]\label{lay:reduced}
For a spherical set $J$, the positive lift of a reduced Coxeter word
has layer amplitude $[0,1]$.  The negative lift, obtained by replacing
each letter by its inverse without reversing the word, has amplitude
$[-1,0]$.  These bounds hold on every pure projective input, and hence
on arbitrary bounded intervals of layers, on either module side.
\end{corollary}

\begin{proof}
First consider a reduced unfolded word.  Its test height can increase
at most once.  Otherwise at a second increase the new letter was a
highest-test vertex of the preceding result, whose height already
exceeded the starting value.  By Lemma~\ref{lay:top-divisor} it can be
extracted on the last-action side of that preceding positive word.
Adding it there would give two equal consecutive Coxeter generators,
contradicting reducedness.

A pure input at $d$ has test height at most $d+2$.  Positive twists
never produce levels below $d$.  If the output had highest layer
$a>d+1$, exterior invariance would force that layer to be supported on
$\widetilde J$.  It has a nonzero top there, so
Lemma~\ref{lay:extremal} gives a nonzero test at $a+2>d+3$.
This contradicts the single-increase bound.  If the initial tests
vanish, they stay zero and give the same contradiction.  This proves
amplitude $[0,1]$ for reduced unfolded words.

By Lemma~\ref{cat:unfolding}, reduced color words unfold to reduced
words, proving the positive assertion.  The inverse of a negative
reduced lift is the positive lift of the reversed reduced expression.
Lemma~\ref{lay:inverse} therefore proves the negative assertion.
The right argument has the same degree bounds.
\end{proof}
